% Part: normal-modal-logic
% Chapter: syntax-and-semantics
% Section: schemas

\documentclass[../../../include/open-logic-section]{subfiles}

\begin{document}

\olfileid[pt]{nml}{syn}{sch}

\olsection{Esquemas e validade}

\begin{defn}
  Um \emph{esquema} é um conjunto de !!{formula}s que contém todas,
  e somente, as instâncias de substituição de alguma !!{formula} modal
  ~$!C$, O.s.\@,
  \[
  \Setabs{!B}{\lexists[!D_1], \dots, \lexists[!D_n] \left(!B =
  \SSubst{!C}{\subst{!D_1}{p_1}, \dots, \subst{!D_n}{p_n}} \right) }.
  \]
  A !!{formula}~$!C$ é chamada !!{formula} \emph{característica} do
  esquema e é única a menos de uma renomeação das !!{propositional
  variable}s. !!^a{formula}~$!A$ é uma \emph{instância} de um esquema
  se pertence ao conjunto.
\end{defn}

É conveniente denotar um esquema pela expressão metalinguística obtida
substituindo `$!A$', `$!B$',~\dots, nos lugares dos componentes atômicos
de~$!C$.
Assim, por exemplo, as expressões a seguir denotam esquemas: `$!A$',
`$!A \lif \Box!A$', `$!A \lif (!B \lif !A)$'. Elas correspondem às
!!{formula}s características $p$, $p \lif \Box p$, $p \lif (q \lif p)$.
O esquema `$!A$' denota o conjunto de \emph{todas} as !!{formula}s.

\begin{defn}
  Um esquema é \emph{verdadeiro} em um modelo se, e somente se, todas
  as suas instâncias são verdadeiras; e um esquema é \emph{válido} se,
  e somente se, é verdadeiro em todos os modelos.
\end{defn}

\begin{prop}\ollabel{prop:Kvalid}
  O esquema \Ax{K} a seguir é válido:
  \begin{equation*}
  \Box(!A \lif !B) \lif (\Box !A \lif \Box !B). \tag{\Ax{K}}
  \end{equation*}
\end{prop}

\begin{proof}
  Precisamos mostrar que todas as instâncias do esquema são verdadeiras
  em todos os mundos de todos os modelos. Sejam, portanto, $\mModel{M} =
  \tuple{W,R,V}$ e $w \in W$ arbitrários. Para mostrar que um condicional
  é verdadeiro em um mundo, supomos que o antecedente seja verdadeiro e
  mostramos que o consequente também é. Neste caso, suponha
  $\mSat{M}{\Box(!A \lif !B)}[w]$ e $\mSat{M}{\Box !A}[w]$. Precisamos
  mostrar $\mSat{M}{\Box !B}[w]$. Seja $w'$ arbitrário tal que $Rww'$.
  Pela primeira hipótese, $\mSat{M}{!A \lif !B}[w']$ e, pela segunda,
  $\mSat{M}{!A}[w']$. Segue-se que $\mSat{M}{!B}[w']$. Como $w'$ era
  arbitrário, $\mSat{M}{\Box !B}[w]$.
\end{proof}

\begin{prop}\ollabel{prop:Dual-valid}
  O esquema \Dual{} a seguir é válido:
  \begin{equation*}
  \Diamond !A \liff \lnot\Box\lnot !A. \tag{\Dual}
  \end{equation*}
\end{prop}

\begin{proof}
  Exercício.
\end{proof}

\begin{prob}
  Demonstre \olref[nml][syn][sch]{prop:Dual-valid}.
\end{prob}

\begin{prop}\ollabel{prop:soundMP}
  Se $!A$ e $!A \lif !B$ são verdadeiras em um mundo de um modelo,
  então $!B$ também é. Logo, as !!{formula}s válidas são fechadas sob
  modus ponens.
\end{prop}

\begin{prop}\ollabel{prop:valid-instances}
  !!^a{formula}~$!A$ é válida sse todas as suas instâncias de substituição
  são válidas. Em outras palavras, um esquema é válido sse sua
  !!{formula} característica é válida.
\end{prop}

\begin{proof}
  A direção ``se'' é evidente, pois $!A$ é uma instância de substituição
  de si mesma.

  Para demonstrar a direção ``somente se'', mostramos o seguinte:
  suponha que $\mModel{M} = \tuple{W, R, V}$ seja um modelo modal e que
  $!B \ident \SSubst{!A}{\subst{!D_1}{p_1},\dots,\subst{!D_n}{p_n}}$
  seja uma instância de substituição de~$!A$. Defina $\mModel{M'} =
  \tuple{W, R, V'}$ por $V'(p_i) = \Setabs{w}{\mSat{M}{!D_i}[w]}$.
  Então $\mSat{M}{!B}[w]$ sse $\mSat{M'}{!A}[w]$, para qualquer~$w \in W$.
  (Deixamos a demonstração como exercício.) Suponha agora que $!A$ seja
  válida, mas que alguma instância de substituição $!B$ de~$!A$ não seja
  válida. Então, para algum $\mModel{M} = \tuple{W, R, V}$ e algum
  $w \in W$, $\mSat/{M}{!B}[w]$. Mas, pela afirmação, temos então
  $\mSat/{M'}{!A}[w]$, e $!A$ não é válida, uma contradição.
\end{proof}

\begin{prob}
  Demonstre a afirmação na parte ``somente se'' da demonstração de
  \olref[nml][syn][sch]{prop:valid-instances}. (Sugestão: use indução
  sobre~$!A$.)
\end{prob}

Observe, porém, que não é verdade que um esquema seja verdadeiro em um
modelo sse sua fórmula característica é verdadeira nele. É claro que a
direção ``somente se'' vale: se toda instância de $!A$ é verdadeira
em~$\mModel{M}$, então a própria $!A$ é verdadeira em~$\mModel{M}$. Mas
pode ocorrer que $!A$ seja verdadeira em~$\mModel{M}$ e que alguma
instância de~$!A$ seja falsa em algum mundo de~$\mModel{M}$. Como um
contraexemplo muito simples, considere~$p$ em um modelo com um único
mundo~$w$ e $V(p) = \{w\}$, de modo que $p$ é verdadeira em~$w$.
Contudo, $\lfalse$ é uma instância de~$p$ e não é verdadeira em~$w$.

\begin{prob}
  Mostre que nenhuma das !!{formula}s a seguir é válida:
  \begin{enumerate}
    \item[\Ax{D}:] \quad $\Box p \lif \Diamond p$;
    \item[\Ax{T}:] \quad $\Box p \lif p$;
    \item[\Ax{B}:] \quad $p \lif \Box\Diamond p$;
    \item[\Ax{4}:] \quad $\Box p \lif \Box \Box p$;
    \item[\Ax{5}:] \quad $\Diamond p \lif \Box \Diamond
      p$.
  \end{enumerate}
\end{prob}

\begin{table}[t]
    \centering
    \begin{tabular}{| l || l |}
      \hline
      {\emph{Esquemas válidos}} & {\emph{Esquemas inválidos}} \\
      \hline\hline
      $\Box(!A \lif !B) \lif (\Diamond !A \lif \Diamond !B)$
      & $\Box (!A \lor !B) \lif (\Box !A \lor \Box !B)$ \\
      $\Diamond (!A \lif !B) \lif (\Box !A \lif \Diamond
      !B)$
      & $(\Diamond !A \land \Diamond !B) \lif \Diamond (!A
      \land !B)$\\
      $\Box (!A \land !B) \liff (\Box !A \land \Box !B)$
      & $!A \lif \Box !A$ \\
      $\Box !A \lif \Box (!B \lif !A)$
      & $\Box \Diamond !A \lif !B$ \\
      $\lnot \Diamond !A \lif \Box (!A \lif !B)$
      & $\Box \Box !A \lif \Box !A$ \\
      $\Diamond (!A \lor !B) \liff (\Diamond !A \lor
      \Diamond !B)$
      & $\Box \Diamond !A \lif \Diamond \Box !A$. \\
      \hline
    \end{tabular}
    \caption{Esquemas válidos e (ou?) inválidos.}
    \ollabel{tab:valid-invalidSchemas}
\end{table}

\begin{prob}%\ollabel{ex:in/validSchemas}
  Demonstre que os esquemas da primeira coluna de
  \olref[nml][syn][sch]{tab:valid-invalidSchemas} são válidos e que os da
  segunda coluna não são válidos.
\end{prob}

\begin{prob}
  Decida se os esquemas a seguir são válidos ou inválidos:
  \begin{enumerate}
  \item $(\Diamond !A \lif \Box !B) \lif (\Box !A \lif \Box
    !B)$;
  \item $\Diamond(!A \lif !B) \lor \Box(!B \lif !A)$.
  \end{enumerate}
\end{prob}

\begin{prob}
  Para cada um dos esquemas a seguir, encontre um modelo $\mModel{M}$
  tal que toda instância da !!{formula} seja verdadeira em $\mModel{M}$:
  \begin{enumerate}
  \item $p \lif \Diamond\Diamond p$;
  \item $\Diamond p \lif \Box p$.
  \end{enumerate}
\end{prob}

\end{document}
